\section{Blind native action census}

The input graph \(\Gsixty\) has sixty vertices, \(120\) edges, and
constant degree four.  Its exact full automorphism action contains
\(480\) permutations.  Before consulting the preserved
\(\Gsixty\)-to-\(\Gthirty\) or \(\Gsixty\)-to-\(\Gfifteen\)
certificates, we enumerated semiregular subgroups of this action.

A subgroup is called semiregular here when every nonidentity element
acts without fixed vertices.  It is cover-admissible when its orbit
projection satisfies the local covering condition and no subgroup
element reverses an edge while fixing that edge setwise.  Subgroups
were then grouped by conjugacy under the full automorphism group.

\begin{proposition}[Blind census]
The native action contains \(198\) semiregular subgroups in \(22\)
full-automorphism conjugacy classes.  All \(22\) classes are
cover-admissible, and none contains an edge-inverting action.
\end{proposition}

The class orders are distributed as follows:
\[
\begin{array}{c|ccccccccc}
\lvert R\rvert &2&3&4&5&6&10&12&20&60\\
\hline
\text{class count}&3&1&4&1&3&3&4&2&1.
\end{array}
\]
Consequently, cover admissibility does not identify a unique native
receipt action.

Three classes are fixed under conjugation by the full automorphism
group.  They have orders \(60\), \(4\), and \(2\).  The order-four
class is a Klein four group, and the order-two class is cyclic.  The
later argument does not choose between these orders.  Instead it asks
conditional questions: if a binary receipt is declared, which binary
action is fully covariant?  If a \(\Vfour\) receipt is declared, which
four-state action is fully covariant?

\begin{remark}
Normality supplies covariance, not receipt-rank selection.  The
existence of normal classes at three different orders is direct
evidence that the raw graph alone has not chosen the apparatus.
\end{remark}
